\begin{abstract}
\noindent
Large language models give a different return forecast each time they are asked. An earlier version of this work proposed to turn this into Black--Litterman inputs: the mean of repeated forecasts as the view and their variance as the view uncertainty $\Omega$, so that the forecaster supplies its own confidence. We re-examine this framework with point-in-time index membership from CRSP, open-weight models tested strictly after their documented training cutoffs, and an out-of-sample choice of the scale parameter $\tau$. We also make the framework's claim testable. Because realized returns contain noise that usually exceeds the view error, calibration must be judged against the predictive variance, and because dispersion grows with volatility, the relevant benchmark is the He--Litterman $\Omega$, not a constant. Two tests were fixed in advance: dispersion is informative about view error beyond volatility, and the dispersion-based $\Omega$ improves the portfolio over conventional choices. For the S\&P~500 in 2025 with gpt-oss-20b, the first test passes (volatility-standardized rank correlation 0.12 within dates, positive on 79\% of dates) and the second fails: the dispersion-based $\Omega$ lowers the Sharpe ratio relative to both the constant and the He--Litterman $\Omega$ under level matching, weight caps, recalibration, alternative risk aversion and transaction costs. Dispersion ranks views roughly correctly but spreads their confidence over two orders of magnitude while their errors differ about fourfold, concentrating the portfolio on views that are not more accurate. With returns as the only input, the views themselves are a transformation of short-term momentum without predictive value. The original performance gains do not survive the corrected design. \pending{update with all markets and models}

\medskip\noindent\textbf{Keywords:} Black--Litterman model; large language models; view uncertainty; forecast calibration; portfolio optimization; look-ahead bias
\end{abstract}
